\section{Introduction}
\label{sec:introduction}
French counting is a clean stress test for multilingual numeracy: it is rule-based, but its lexical surface is atypical in the \irregularband range. Errors in numeric language matter in tutoring, translation, and safety-critical communication. Our main question is simple: do strong LLMs degrade on French irregular forms, or can they map these forms consistently?

Prior numeracy studies establish that tokenization, representation choice, and synthetic pretraining all affect number handling \citep{wallace2019nlp,geva2020injecting,jin2021numgpt}. Multilingual reasoning benchmarks such as MGSM report aggregate French performance \citep{shi2022language}, and interpretability work identifies arithmetic-sensitive components \citep{yu2024interpreting,chang2024unraveling}. However, existing evaluations do not isolate French-specific numeral morphology with full coverage and band-level significance tests.

We address this gap with \frcountmap, an inference-only probe suite for French numerals. We evaluate deterministic API outputs on \digitword, \worddigit, and \patterncls tasks over a synthetic 0--199 dataset, then connect those results to \mgsmfr performance. Our quantitative preview is direct: transduction is near-perfect and does not show a significant irregular-band penalty, while \patterncls exhibits large and significant irregular-band advantages under our label design.

In summary, our main contributions are:
\begin{itemize}
    \item We introduce a reproducible French numeral probe pipeline with complete 0--199 coverage and transparent normalization/scoring.
    \item We conduct controlled model-by-condition comparisons showing ceiling-level \worddigit and near-ceiling \digitword performance without a robust irregular-range drop.
    \item We provide band-wise risk-difference tests with FDR correction and a pooled logistic model that quantifies where band effects do and do not appear.
    \item We show that prompt reasoning affects broader French math reasoning (\mgsmfr) more than numeral transduction itself.
\end{itemize}

\para{Paper organization.} \Secref{sec:related_work} discusses prior numeracy and multilingual reasoning research. \Secref{sec:methodology} details data, probes, and metrics. \Secref{sec:results} presents quantitative findings and statistical tests. \Secref{sec:discussion} covers interpretation, limitations, and implications, and \secref{sec:conclusion} concludes.
